\chapter{Termo de Abertura do Projeto}

\section{Dados do projeto}
\begin{description}
    \item [Nome do Projeto:] MrBombastic
    \item [Data de abertura:] 09/04/2026
    \item [Código:] 1-A
    \item [Patrocinador:] Universidade de Brasília
    % \item [Responsável:] ???
    \item [Gerente:] Tiago Lemes Teixeira / 231026581 / 231026581@aluno.unb.br / (61) 99584-1973
\end{description}

\section{Objetivos}

\textcolor{black}{
O objetivo do projeto é desenvolver um micromouse autônomo capaz de navegar de forma eficiente em labirintos de diferentes dimensões, alcançando a área central do labirinto sem intervenção humana após sua ativação inicial. O sistema deve realizar detecção de paredes, mapeamento progressivo do ambiente, monitoramento contínuo da posição, tomada de decisões de navegação e registro de telemetria em tempo real. Este objetivo atende aos critérios SMART da seguinte forma:
}

\begin{description}
    \item[\textit{Specific} (específico):] Desenvolver um micromouse autônomo capaz de interpretar o ambiente, identificar paredes, mapear o labirinto e encontrar o caminho até a área central, operando de forma completamente independente. O projeto também inclui a implementação de um sistema de telemetria para registro de trajetória, velocidade, consumo de bateria e métricas de desempenho.
    
    \item[\textit{Measurable} (mensurável):] O desempenho será avaliado pela capacidade do micromouse de completar com sucesso três configurações de labirinto (4x4, 8x8 e 16x16), registrando trajetória percorrida, tempo total, velocidade média e consistência das decisões de navegação.
    
    \item[\textit{Agreed} (acordado):] O objetivo foi definido em conjunto pela equipe do projeto, sob supervisão da Universidade de Brasília, considerando as demandas das áreas de eletrônica, energia, estrutura e software, estando alinhado com as expectativas acadêmicas e formativas das partes envolvidas.
    
    \item[\textit{Realistic} (realista):] O desenvolvimento do micromouse é viável diante dos recursos técnicos disponíveis, do orçamento destinado e das competências multidisciplinares da equipe (sensores, controle, eletrônica e programação), compatíveis com o escopo acadêmico do projeto.
    
    \item[\textit{Time Bound} (limitado no tempo):] O projeto deverá ser concluído até o final do semestre 2026.1, com marcos intermediários estabelecidos no cronograma acadêmico.
\end{description}

\section{Mercado-alvo}

\textcolor{black}{
O mercado-alvo do projeto Micromouse Autônomo engloba diferentes públicos que podem se beneficiar direta ou indiretamente dos resultados e soluções desenvolvidas. Entre eles:
}

\begin{itemize}
    \item \textbf{Estudantes e Pesquisadores de Engenharia:} \\
    O projeto funciona como uma plataforma prática para o estudo de robótica móvel, sensores, algoritmos de navegação, mapeamento e sistemas embarcados, fortalecendo o aprendizado.

    \item \textbf{Competições e Comunidades de Robótica:} \\
    O micromouse é amplamente utilizado em torneios e eventos de tecnologia, onde equipes testam soluções de navegação autônoma. Assim, o projeto contribui para a formação de equipes competitivas e para o avanço de práticas de inovação em ambientes estudantis e científicos.

    \item \textbf{Comunidade Acadêmica e Científica:} \\
    De forma indireta, o projeto promove o avanço da pesquisa em robótica de pequena escala, telemetria, otimização de rotas e sistemas autônomos de baixo custo, incentivando a formação de novos profissionais na área tecnológica.

\end{itemize}

\section{Requisitos}

Os requisitos do MrBombastic definem as funcionalidades e restrições essenciais que o sistema deve cumprir para garantir seu propósito educacional e tecnológico. O sistema neste contexto de requisitos refere-se ao conjunto de algoritmos e lógica computacional. Esses requisitos estão organizados em Requisitos Funcionais (RF) e Requisitos Não Funcionais (RNF), que abrangem, respectivamente, os comportamentos do sistema e as limitações técnicas necessárias para assegurar desempenho, confiabilidade e padronização do projeto.

\begin{table}[htbp]
\begin{center}
\scriptsize
\setlength{\tabcolsep}{4pt}
\renewcommand{\arraystretch}{1}
\caption{Requisitos Funcionais do Sistema (Integrado G2).}
\begin{tabular}{|c|p{13cm}|}
\hline
\textbf{ID} & \textbf{Descrição} \\
\hline
\textbf{RF01} & O sistema deve permitir que o micromouse navegue de forma autônoma em labirintos sem intervenção humana após a ativação inicial. \\ \hline
\textbf{RF02} & O sistema deve ser capaz de detectar paredes durante a navegação. \\ \hline
\textbf{RF03} & O sistema deve monitorar continuamente a posição do micromouse. \\ \hline
\textbf{RF04} & O sistema deve mapear progressivamente o ambiente do labirinto. \\ \hline
\textbf{RF05} & O sistema deve identificar quando o micromouse alcança a área central do labirinto. \\ \hline
\textbf{RF06} & O sistema deve interpretar o ambiente com base na disposição das paredes. \\ \hline
\textbf{RF07} & O sistema deve tomar decisões de navegação com base nas rotas disponíveis. \\ \hline
\textbf{RF08} & O sistema deve adaptar a navegação a diferentes dimensões de labirinto (4x4, 8x8 e 16x16). \\ \hline
\textbf{RF09} & O sistema deve registrar a trajetória percorrida pelo micromouse. \\ \hline
\textbf{RF10} & O sistema deve calcular o consumo de bateria. \\ \hline
\textbf{RF11} & O sistema deve exibir o consumo de bateria. \\ \hline
\textbf{RF12} & O sistema deve calcular a velocidade média do micromouse. \\ \hline
\textbf{RF13} & O sistema deve exibir a velocidade média do micromouse. \\ \hline
\textbf{RF14} & O sistema deve calcular o tempo total de execução. \\ \hline
\textbf{RF15} & O sistema deve exibir o tempo total de execução. \\ \hline
\textbf{RF16} & O sistema deve exibir em tempo real os dados de telemetria por meio de uma interface web. \\ \hline
\textbf{RF17} & O sistema deve exibir o tipo de labirinto. \\ \hline
\textbf{RF18} & O sistema deve exibir o trajeto percorrido. \\ \hline
\textbf{RF19} & O sistema deve indicar o status do desafio (cumprido ou não). \\ \hline
\textbf{RF20} & O sistema deve armazenar os dados das execuções em um banco de dados após a conclusão. \\ \hline
\textbf{RF21} & O sistema deve permitir a consulta dos dados armazenados por labirinto. \\ \hline
\textbf{RF22} & O sistema deve permitir a comparação de desempenho entre diferentes execuções. \\ \hline
\textbf{RF23} & O sistema deve ser capaz de achar o caminho mais rápido após mapeamento do labirinto. \\ \hline
\textbf{RF24} & O sistema deve exibir o labirinto percorrido após a execução. \\ \hline
\end{tabular}
\end{center}
\label{RF}
\end{table}

\begin{table}[htbp]
\begin{center}
\scriptsize
\setlength{\tabcolsep}{4pt}
\renewcommand{\arraystretch}{1}
\caption{Requisitos Não Funcionais do Sistema (Integrado G2).}
\begin{tabular}{|c|p{10cm}|c|}
\hline
\textbf{ID} & \textbf{Descrição} & \textbf{Categoria URPS+} \\
\hline
\textbf{RNF01} & O sistema deve se comunicar com o sistema por wi-fi. & Suportabilidade \\ \hline
\textbf{RNF02} & O micromouse deve parar completamente no tempo de 2s, após receber comando de parada emergência. & Confiabilidade \\ \hline
\textbf{RNF03} & O sistema deve ser capaz de armazenar dados no banco de dados relacional de forma autônoma. & Confiabilidade \\ \hline
\textbf{RNF04} & O micromouse deve ser capaz de fazer uma decisão ao chegar uma célula no tempo de 5s. & Restrição \\ \hline
\textbf{RNF05} & O sistema de interface web deve exibir dados atualizados do micromouse de telemetria sem exceder 500 ms. & Desempenho \\ \hline
\textbf{RNF06} & O sistema de alimentação (bateria) não deve exceder 25\% do peso total do micromouse. & Restrição \\ \hline
\end{tabular}
\end{center}
\label{RNF}
\end{table}

\section{Justificativa}

\textcolor{black}{
Por se tratar de um desafio que integra sensores, atuadores, programação de estratégias de busca e otimização de rotas, o projeto oferece aos estudantes uma oportunidade significativa para preencher lacunas da formação acadêmica, permitindo a aplicação direta de conceitos teóricos em um sistema físico real. Dessa forma, o micromouse serve como plataforma educacional robusta, contribuindo para o desenvolvimento de habilidades essenciais de engenharia e para a compreensão das aplicações práticas dessa tecnologia no contexto profissional.
}

\section{Indicadores}
\textcolor{black}{Alguns dos indicadores que determinam o mercado consumidor e o potencial de aplicação do \textbf{MrBombastic} são:}
\begin{enumerate}
    \item As competições de Micromouse são realizadas há quase 50 anos. O primeiro anúncio de um evento do tipo foi feito pela revista IEEE Spectrum em 1977, que apresentou a proposta do “Amazing Micro-Mouse Maze Contest”, e a primeira competição oficial ocorreu em Nova York em 1979 \cite{christiansen1977}.
    \item O interesse por competições de Micromouse no Brasil vem crescendo, e um exemplo recente é a 3ª edição do RoboChallenge Brasil, realizada pelo Instituto Mauá de Tecnologia em 2024, que contou com mais de 200 robôs competindo em diversas categorias, incluindo Micromouse \cite{maua_robochallenge_2024}.
\end{enumerate}

% \textcolor{red}{Listar até 10 indicadores que determinam o mercado consumidor do produto desenvolvido: exemplo: 1) n° de alunos da FGA que utilizam ônibus às 18:00; 2) n° de usuários do restaurante universitários, 3) número de idosos classificados como público-alvo no DF e no estado de Goiás, 4) n° de empresas de segurança registradas no DF etc.}
% \section{Aprovações}

% \begin{tabular}{ l l }
%   \textbf{Patrocinador:} & \_\_\_\_\_\_\_\_\_\_\_\_\_\_\_\_\_\_\_\_\_\_\_\_\_\_\_\_\_\_\_\_\_\_ \\
%   & \\
%   \textbf{CEO:} & \_\_\_\_\_\_\_\_\_\_\_\_\_\_\_\_\_\_\_\_\_\_\_\_\_\_\_\_\_\_\_\_\_\_ \\
%   & \\
%   \textbf{Gerente:} & \_\_\_\_\_\_\_\_\_\_\_\_\_\_\_\_\_\_\_\_\_\_\_\_\_\_\_\_\_\_\_\_\_\_ \\
% \end{tabular}
